\section*{序論}
10月26日朝までの24時間は、大型モデルの公式発表よりも、オープンウェイトモデルの実用性や既存エージェントの利用体験を巡る議論が目立った。Nanbeige4.3-33Bの公開は、約30B級の密モデルをエージェントに使う際の性能と必要メモリに関心を集めた。Claude Duetについては、プロアクティブな作業を担う研究プレビューの利用報告が続いている。

一方、AnthropicのIPOやGPT-6.1 Astraの提供時期については、未確認の投稿が話題となった。任天堂の宮本茂氏による生成AIへの発言とされる引用も、原インタビューの確認が限定的なまま共有されている。こうした話題では、X上で観測された主張と確定した発表を分けて読む必要がある。

10月のモデル動向を振り返るダイジェストや、AI生成広告の表示規制に関する再言及も見られた。本号はGrokが収集した8件を記載順に紹介する。観測窓の終端を越えた時刻の投稿は、その旨を明記した。
